% Part: first-order-logic
% Chapter: completeness
% Section: downward-ls

\documentclass[../../../include/open-logic-section]{subfiles}

\begin{document}

\olfileid{fol}{com}{dls}
\olsection{مبرهنة لوفنهايم--سكولم}

تقول مبرهنة لوفنهايم--سكولم إنه إذا كان لنظرية نموذج غير منتهٍ، فلها
أيضًا نموذج !!{denumerable} على الأكثر. ومن النتائج المباشرة لهذه
الحقيقة أن منطق الرتبة الأولى لا يستطيع التعبير عن كون حجم
!!a{structure} غير !!{enumerable}: فكل !!{sentence} أو مجموعة من
ال!!{sentence}s تصدق في جميع ال!!{structure}s غير ال!!{enumerable} تصدق
أيضًا في بنية !!{enumerable} ما.

\begin{thm}
\ollabel{thm:downward-ls} إذا كانت $\Gamma$ متسقة، فلها نموذج
!!a{enumerable}، أي إنها قابلة للإرضاء في !!a{structure} مجالها إما
منتهٍ وإما !!{denumerable}.
\end{thm}

\begin{proof}
إذا كانت $\Gamma$ متسقة في اللغة~$\Lang L$، فإن برهان مبرهنة الاكتمال
يبني امتداد هنكني~$\Lang{L'}$ وبنية حدود عليه لا يزيد مجالها على مجموعة
الحدود المغلقة~$\Trm[L']_0$ في~$\Lang{L'}$. ومن ثم تكون هذه البنية
!!{denumerable} على الأكثر، ويكون ردّها~$\Struct M$ إلى اللغة
الأصلية~$\Lang L$ نموذجًا لـ~$\Gamma$ له المجال نفسه.
\end{proof}

\begin{thm}
\ollabel{noidentity-ls} إذا كانت $\Gamma$ مجموعة متسقة من ال!!{sentence}s
في لغة منطق الرتبة الأولى الخالية من ال!!{identity}، فلها نموذج
!!a{denumerable}، أي إنها قابلة للإرضاء في !!a{structure} مجالها غير
منتهٍ و!!{enumerable}.
\end{thm}

\begin{proof}
إذا كانت $\Gamma$ متسقة ولا تحتوي أي !!{sentence} ترد فيها ال!!{identity}،
فإن بنية الحدود في لغة هنكن~$\Lang L'$ التي يبنيها برهان مبرهنة
الاكتمال لها مجال~$\Domain{M}$ مطابق لمجموعة الحدود المغلقة في تلك
اللغة. ومن ثم يكون ردّ~$\Struct{M}$ إلى اللغة الأصلية نموذجًا
!!{denumerable} لـ~$\Gamma$، لأن $\Trm[L']_0$ !!{denumerable}.
\end{proof}

\begin{ex}[مفارقة سكولم]
نظرية المجموعات لزرميلو--فرينكل~$\Log{ZFC}$ إطار بالغ القوة يمكن فيه
التعبير عن جميع العبارات الرياضية تقريبًا، ومنها الحقائق المتعلقة
بأحجام المجموعات. فمثلًا تستطيع $\Log{ZFC}$ أن تبرهن أن مجموعة الأعداد
الحقيقية~$\Real$ غير !!{enumerable}، وأن تبرهن مبرهنة كانتور القائلة إن
مجموعة أجزاء أي مجموعة أكبر من المجموعة نفسها، وهلم جرًا. وإذا كانت
$\Log{ZFC}$ متسقة، كانت جميع نماذجها غير منتهية، وفوق ذلك تحتوي كلها
!!{element}s تقول النظرية عنها إنها غير !!{enumerable}، مثل العنصر الذي
يحقق مبرهنة~$\Log{ZFC}$ القائلة بوجود مجموعة أجزاء الأعداد الطبيعية.
وبحسب مبرهنة لوفنهايم--سكولم، لـ~$\Log{ZFC}$ أيضًا نماذج
!!{enumerable}، أي نماذج تحتوي مجموعات «غير !!{enumerable}» لكنها هي
نفسها !!{enumerable}.
\end{ex}

\end{document}
